\section*{Bab \ref{ch:b1:organic-redox} --- Oksidasi dan Reduksi dalam Kimia Organik}

\begin{solution}{exo:b1:organic-redox:1}
Etanol: \ce{CH3} $-3$, \ce{CH2OH} $-1$. Etanal: \ce{CH3} $-3$, CHO $+1$.
Asam etanoat: \ce{CH3} $-3$, COOH $+3$. Propanon: \ce{CH3} $-3$ (dua
kali), C=O $+2$. Asam metanoat: $+2$.
\end{solution}

\begin{solution}{exo:b1:organic-redox:2}
Propan-1-ol: propanal, lalu asam propanoat. Propan-2-ol: propanon.
2-Metilpropan-2-ol: tidak bereaksi (warna jingganya tetap).
\end{solution}

\begin{solution}{exo:b1:organic-redox:3}
\ce{CH3COCH3 + 2H+ + 2e- -> CH3CH(OH)CH3};
\ce{CH3CHO + 2H+ + 2e- -> CH3CH2OH}.
\end{solution}

\begin{solution}{exo:b1:organic-redox:4}
\ce{NaBH4}: butan-1-ol, butan-2-ol, tidak bereaksi, tidak bereaksi (asam
itu hanya dideprotonasi). \ce{LiAlH4}: butan-1-ol, butan-2-ol, butan-1-ol
dan etanol, butan-1-ol.
\end{solution}

\begin{solution}{exo:b1:organic-redox:5}
\ce{5CH3CH(OH)CH3 + 2MnO4- + 6H+ -> 5CH3COCH3 + 2Mn^2+ + 8H2O}.
$n = 1.20/60.0 = \qty{0.0200}{mol}$; permanganat $\frac25 \times 0.0200 =
\qty{8.0e-3}{mol}$, \qty{400}{mL} larutan \qty{0.020}{mol/L}.
\end{solution}

\begin{solution}{exo:b1:organic-redox:6}
Panaskan alkohol dengan \omterm{def:b1:nernst:couple}{oksidatornya} dalam labu yang dipasang untuk
distilasi, pada suhu di antara kedua titik didih: propanal
(\qty{48}{\celsius}) terdistilasi segera setelah terbentuk dan luput dari
oksidasi lebih lanjut; propan-1-ol (\qty{97}{\celsius}) tetap di dalam
labu.
\end{solution}

\begin{solution}{exo:b1:organic-redox:7}
Sepasang elektron B--H dari \ce{BH4-} menyerang karbon karbonil, sementara
pasangan $\pi$ C=O berpindah ke oksigen; \omterm{def:b1:alcohol-activation:alkoxide}{alkoksidanya} mengambil sebuah
proton dari etanol. Hidrogen pada karbon produk (\ce{CH3CH2OH}, salah satu
dari dua hidrogen pada C1) berasal dari pereaksi; hidrogen pada oksigen
berasal dari pelarut.
\end{solution}

\begin{solution}{exo:b1:organic-redox:8}
Butan-2-ol \omterm{def:b1:stereochemistry-in-depth:stereocentre}{kiral}, tetapi keton yang datar diserang sama banyak pada kedua
sisinya: \omterm{def:b1:stereochemistry-in-depth:enantiomers}{campuran rasemat}, tidak aktif optis.
\end{solution}

\begin{solution}{exo:b1:organic-redox:9}
\ce{3CH3CH2OH + 2Cr2O7^2- + 16H+ -> 3CH3COOH + 4Cr^3+ + 11H2O}: dikromat
yang jingga (kromium $+$VI) menjadi kromium(III) yang hijau.
\end{solution}

\begin{solution}{exo:b1:organic-redox:10}
2-Metilbutan-2-ol: etanal + \ce{CH3CH2MgBr} menghasilkan butan-2-ol
(sebuah adisi, bukan redoks); oksidasi menjadi butanon; \ce{CH3MgBr} pada
butanon, lalu pengolahan akhir. Satu oksidasi. Pentan-3-ol: propanal +
\ce{CH3CH2MgBr} dalam satu tahap, tanpa redoks.
\end{solution}

\begin{solution}{exo:b1:organic-redox:11}
1. \omterm{def:b1:acetals-protection:protecting-group}{Proteksi} keton sebagai \omterm{def:b1:acetals-protection:hemiacetal}{asetal} siklik (etana-1,2-diol, asam,
Dean--Stark). 2. \ce{LiAlH4} dalam eter kering mereduksi ester menjadi
alkohol primer (dan metanol). 3. Asam berair melepaskan \omterm{def:b1:acetals-protection:hemiacetal}{asetalnya}:
5-hidroksipentan-2-on, \ce{CH3CO(CH2)3OH}.
\end{solution}

\begin{solution}{exo:b1:organic-redox:12}
\ce{RCH2OH}: karbon karbinol pada $-1$; \ce{RCOOH}: $+3$; empat elektron,
seperti dalam \ce{RCH2OH + H2O -> RCOOH + 4H+ + 4e-}. Metana ($-4$)
menjadi karbon dioksida ($+4$): delapan elektron,
\ce{CH4 + 2H2O -> CO2 + 8H+ + 8e-}.
\end{solution}

\begin{solution}{pb:b1:organic-redox:1}
\textbf{1.} C1 (yang membawa OH): 0; kelima \ce{CH2}: $-2$.
\textbf{2.} C1 (C=O): $+2$; kelima \ce{CH2}: $-2$.
\textbf{3.} Kedua karbon COOH: $+3$; keempat \ce{CH2}: $-2$.
\textbf{4.} C1 ($0 \to +3$) dan tetangganya C6 ($-2 \to +3$), yang menjadi
karbon karboksilat kedua; ikatan C1--C6 putus.
\textbf{5.} $3 + 5 = 8$ elektron.
\textbf{6.} Asam nitrat cukup kuat untuk memutus ikatan C--C di sebelah
karbonil, yang tidak dilakukan oleh \omterm{def:b1:nernst:couple}{oksidator} yang lebih lunak.
\textbf{7.} \ce{C6H12O + 3H2O -> C6H10O4 + 8H+ + 8e-}.
\textbf{8.} $+$V dan $+$I.
\textbf{9.} \ce{2HNO3 + 8H+ + 8e- -> N2O + 5H2O}.
\textbf{10.} Jika dijumlahkan, delapan elektron, \ce{8H+}, dan tiga
molekul air saling meniadakan:
\ce{C6H12O + 2HNO3 -> C6H10O4 + N2O + 2H2O}.
\textbf{11.} 8.
\textbf{12.} $10^6/146.0 = \qty{6.85e3}{mol}$ asam; $2 \times 6.85 \times
10^3 \times 63.0 = \qty{8.6e5}{g}$, \qty{0.86}{t} asam nitrat.
\textbf{13.} \ce{H2O2 + 2H+ + 2e- -> 2H2O}.
\textbf{14.} \ce{C6H12O + 4H2O2 -> C6H10O4 + 5H2O}.
\textbf{15.} $4 \times 6.85 \times 10^3 \times 34.0 = \qty{9.3e5}{g}$,
\qty{0.93}{t}.
\textbf{16.} Air: tidak ada oksida nitrogen.
\textbf{17.} Tetapan kesetimbangan yang besar tidak mengatakan apa pun
tentang laju; \ce{H2O2} bereaksi lambat tanpa \omterm{def:b1:elementary-steps:catalyst}{katalis}, seperti halnya
penguraiannya sendiri.
\textbf{18.} Natrium borohidrida mereduksi sikloheksanon kembali menjadi
sikloheksanol: \ce{4C6H10O + BH4- -> B(OC6H11)4-}, lalu
\ce{B(OC6H11)4- + 4H2O -> 4C6H11OH + B(OH)4-}.
\textbf{19.} \qty{6.85e3}{mol}.
\textbf{20.} \qty{6.85e3}{mol} \ce{N2O} (satu per molekul asam).
\textbf{21.} $6.85 \times 10^3/0.95 \times 100.0 = \qty{7.2e5}{g}$, \qty{0.72}{t}.
\textbf{22.} $6.85 \times 10^3 \times 44.0 = \qty{3.0e5}{g}$: \textbf{\qty{0.30}{t}
\ce{N2O} per ton asam adipat}.
\end{solution}
